%% delta_betting.tex
%% Sections: Betting Markets, Welfare Implications, Conclusion
%% Paper: "The Economics of Extended Play"
%% ---------------------------------------------------------

\section{Betting Markets: Pricing a Structural Regime Change}\label{sec:betting}

Sports betting markets offer a high-frequency laboratory for studying how
prices respond to regulatory change. The FIFA directive altered the temporal
structure of matches---extending second-half stoppage time by roughly seven
minutes on average---but left the fundamental object of most wagers, total
goals, essentially unchanged. This section asks whether sophisticated and
retail betting markets distinguished between these two facts.

\subsection{Overall Forecast Accuracy}\label{subsec:brier}

We evaluate forecast accuracy using Brier scores for match-outcome (1X2)
and over/under markets, estimating difference-in-differences specifications
of the form used throughout the paper. If the directive degraded
bookmakers' ability to forecast outcomes, we would observe a positive shift
in Brier scores (worse calibration) in the post-directive period for
treated leagues relative to controls.

Table~\ref{tab:brier} reports the results. For Pinnacle, widely regarded
as the sharpest global bookmaker, the Brier-score DiD estimate is $+0.007$
(SE~$= 0.008$, $p = 0.41$, $N = 27{,}867$). For Bet365, the
largest retail sportsbook by volume, the estimate is $+0.005$
(SE~$= 0.008$, $p = 0.58$, $N = 28{,}382$). Estimates for the
remaining bookmakers in our sample---Betway, Interwetten, William Hill,
and bet-at-home---are of similar magnitude and similarly insignificant.
Across the full panel, there is no evidence that the directive impaired
forecast accuracy for any bookmaker.

Two non-competing explanations are consistent with this null result.
First, markets may have rapidly incorporated the regime change, adjusting
lines within the opening weeks of the post-directive period such that
forecast accuracy was restored before it could register in
season-level comparisons. Second---and more parsimoniously---the directive
may not have changed the object being forecasted. Since total goals per
match remained stable (Section~\ref{sec:goals}), the over/under lines that
bookmakers set before kickoff required little revision: the same number of
goals was scored, merely redistributed across the ninety-plus minutes.
Both explanations imply well-functioning markets; they differ only in
whether there was anything to which markets needed to adjust.

\subsection{Bookmaker Disagreement and Market Segmentation}\label{subsec:disagreement}

A standard prediction of models with heterogeneous information is that
less-informed agents adjust more slowly to structural breaks, producing
transient disagreement between ``sharp'' (professional, low-margin) and
``soft'' (retail, higher-margin) bookmakers \citep{Levitt2004}. If the
directive constituted a genuine informational shock, we would expect
Pinnacle's forecasts to dominate Bet365's in the immediate
post-directive period, with the gap narrowing over time.

We find no evidence of such a pattern. The difference in Brier-score DiD
between the sharp-bookmaker average and the retail-bookmaker average is
economically negligible and statistically insignificant. Both market
segments absorbed the regime change---or its irrelevance to total
goals---symmetrically. This null is informative: it suggests that the
directive did not generate the kind of private-information advantage that
separates sharp from retail books in other settings, such as injury news
or lineup changes \citep{Croxson2014}.

\subsection{Speed of Adjustment}\label{subsec:halflife}

To characterize the dynamics of any market response, we estimate the
``half-life'' of excess miscalibration following the directive. We fit
exponential-decay models to the week-by-week Brier-score residuals in
both the 1X2 and over/under markets, asking how quickly any initial
deterioration in accuracy dissipated.

The exercise yields a stark result: for the 1X2 market, the exponential
fit fails entirely because there is no detectable initial excess to
decay. For the over/under market, the same pattern holds---there is no
positive excess in the opening weeks from which a decay trajectory could
be estimated. The most straightforward interpretation is that markets
adjusted so quickly that no measurable adjustment period existed at the
matchweek frequency, or---consistent with the parsimonious explanation
offered above---that there was no adjustment to make.

From the standpoint of the efficient markets hypothesis, this finding is
perhaps unsurprising. The directive was publicly announced, universally
applied, and easily observable in real time. It lacked the ambiguity that
characterizes settings where markets have been shown to adjust slowly,
such as changes in player ability \citep{Massey2013} or the
incorporation of complex statistical signals \citep{Brown1993}. The
directive changed \emph{when} goals occurred but not \emph{how many}---and
over/under lines, the most directly affected product, showed no
systematic miscalibration at any horizon.

\subsection{Closing-Line Value}\label{subsec:clv}

Closing-line value (CLV) measures whether late-arriving information
systematically moves odds in one direction, providing a test of whether
informed bettors exploited knowledge of the directive's implications
\citep{Kaunitz2017}. We regress the closing-line drift (the change in
implied probability between opening and closing lines) on a
post-directive indicator, controlling for league and season fixed effects.

For the 1X2 home-win market, the post-directive coefficient is $0.021$
($p = 0.092$). For the over/under ``over'' market, the coefficient is
$0.006$ ($p = 0.082$). Neither estimate is significant at conventional
levels, though both are marginally suggestive of slight late-market
learning---a pattern consistent with a small number of sophisticated
bettors incorporating match-specific stoppage-time expectations into
their models. Notably, the Pinnacle over/under closing-line drift
coefficient is $0.002$ ($p = 0.60$), indicating that the sharpest
benchmark in our sample exhibited no systematic drift whatsoever. The
marginal effects, where they appear, are confined to softer books and are
too small to survive standard multiple-testing corrections.

\subsection{First-Half Goals as Placebo}\label{subsec:placebo_betting}

As an additional specification check, we estimate the directive's effect
on first-half goals. Since the directive was explicitly targeted at
second-half stoppage time, first-half goal totals should be unaffected by
any mechanism through which the directive operates. This placebo test
yields a coefficient of $\beta = -0.023$ ($p = 0.33$), confirming that
the treatment is specific to second-half timing rather than a general
shift in match structure or refereeing behavior affecting goal
production.\footnote{We note, however, that first-half stoppage time
itself increased by 0.550 minutes ($p < 0.001$), suggesting that the
directive---or the accompanying cultural shift in timekeeping---spilled
over into first-half officiating. The placebo thus passes for goals but
fails for stoppage time, a distinction we discuss further in
Section~\ref{sec:mechanisms}.}

\subsection{Discussion}

The betting-market evidence supports a consistent narrative: the FIFA
directive did not generate exploitable inefficiencies in the pricing of
match outcomes or goal totals. Markets showed no deterioration in
accuracy, no differential adjustment between sharp and retail segments, no
measurable adjustment period, and no economically significant closing-line
drift. The most parsimonious interpretation is that betting markets had
little \emph{to} adjust. The directive changed the temporal distribution
of goals but not their total, and the products most directly affected by
this redistribution---over/under and correct-score markets---were priced
off the same underlying total that remained stable. This finding
complements the goal-production results in Section~\ref{sec:goals} and
underscores a broader point: regulatory interventions that alter the
\emph{structure} of a market without altering its \emph{fundamentals} may
pass through the price system without leaving a trace.


%% ---------------------------------------------------------

\section{Welfare Implications}\label{sec:welfare}

The preceding sections have established the directive's effects on match
outcomes, scoring patterns, and market prices. We now turn to the
normative question: who gained, who lost, and on net, was the directive
welfare-improving? We examine four constituencies: consumers (fans and
viewers), workers (players), clubs (through competitive balance), and
broadcasters.

\subsection{Consumer Welfare: The Redistribution of Entertainment}\label{subsec:consumer}

From the perspective of match-day entertainment, the directive's most
visible consequence was a dramatic increase in late-game drama. Goals
scored after the 90th minute rose by approximately 20 percent, and the
probability of a meaningful goal in added time---defined as a goal after
the 85th minute in a match decided by one goal or fewer---increased by a
similar magnitude. The probability that stoppage time exceeded seven
minutes rose from 14.8 percent to 39.2 percent, transforming the final
minutes from a brief coda into a substantive phase of play.

Yet total goals per match were unchanged. The entertainment ``budget,'' if
measured in goals, was redistributed rather than expanded: more late goals,
but commensurately fewer goals per minute during the expanded playing
time. Whether this redistribution increased consumer surplus depends on
the shape of the fan's utility function over goal timing. If fans exhibit
a preference for dramatic, late-deciding goals---as the behavioral
literature on peak-end effects would suggest \citep{Kahneman1993}---then
the directive improved the quality of the viewing experience even while
holding quantity fixed. The disproportionate media attention paid to
stoppage-time goals, from Sergio Agüero's 2012 title-winner to the
routine last-gasp equalisers that dominate highlight reels, provides
circumstantial evidence for such a preference.

From a broadcasting perspective, the directive increased the duration of
live match coverage by approximately seven minutes per match. This
expansion creates additional advertising inventory and, more importantly,
extends the window during which audiences remain engaged. Whether this
translates into higher rights fees is an empirical question we cannot
resolve without viewership data, but the revealed preferences of
broadcasters---who have not objected to the longer matches---are
suggestive.

\subsection{Worker Welfare: Player Health and Workload}\label{subsec:worker}

The directive's most direct labor-market consequence was an increase in
per-match workload. Our player-level estimates indicate an increase of
$0.76$ minutes per match per player (SE~$= 0.145$, $p < 0.001$,
$N = 830{,}569$ player-match observations). Over a season, this
accumulates to a non-trivial increase in total playing time, raising
legitimate concerns about fatigue and injury risk---concerns that have
been voiced prominently by players' unions and medical staff
\citep{Ekstrand2023}.

The injury evidence, however, tells a more nuanced story. Total injuries
per player-season \emph{decreased} following the directive ($\beta =
-0.090$, SE~$= 0.032$, $p = 0.005$; pre-treatment mean $= 1.886$). This
counterintuitive finding admits several explanations. First, the
reduction in per-minute intensity documented in Section~\ref{sec:goals}
may have lowered the biomechanical stress per unit of playing time. If
the same total effort is distributed over more minutes, the peak
physiological loads that trigger acute injuries may be attenuated.
Second, the longer stoppage periods may have provided de facto recovery
windows---additional time for substitutions, hydration, and tactical
reorganization---that functioned as in-match rest. Third, the finding may
reflect confounding changes in squad rotation or fixture scheduling that
coincided with the directive.

The most policy-relevant test concerns muscular injuries specifically,
since these are the injury type most directly linked to workload and
fatigue \citep{Ekstrand2016}. The DiD estimate for muscular injuries is
$\beta = -0.007$ (SE~$= 0.019$, $p = 0.73$; pre-treatment mean $=
0.525$). This null is well-powered: our design has 88.6 percent power to
detect a 10 percent increase in muscular injury rates, based on
$N = 23{,}255$ player-seasons. The data are therefore informative enough
to reject a clinically meaningful increase, and we conclude that there is
no evidence that the directive harmed player welfare through increased
muscular injury risk.

We emphasize appropriate caution in interpreting these results. The injury
decrease could reflect time-varying confounders---improvements in sports
medicine, changes in training protocols, or strategic rest
patterns---that coincided with but were not caused by the directive.
Our identification strategy absorbs league and season fixed effects but
cannot rule out within-league temporal shocks correlated with the
directive's adoption. The most defensible claim is that the directive did
not \emph{measurably worsen} the injury environment, not that it improved
it.

\subsection{Competitive Balance}\label{subsec:balance}

A recurring concern about rule changes in professional sports is that
they may systematically advantage certain types of clubs---wealthier teams
with deeper squads, for instance, or defensive teams that benefit from
running down the clock \citep{Fort2007}. We test for such effects by
examining the standard deviation of final league points and the normalized
Herfindahl-Hirschman Index (NHHI) of wins in our treated leagues.

Neither measure shows a significant change following the directive. The
redistribution of scoring opportunities across the match---more goals in
added time, commensurately fewer in regulation---did not systematically
advantage any team type. This null is consistent with the aggregate
goal-production results: if total goals per match are unchanged, and the
within-match redistribution is symmetric across home and away teams, there
is no mechanical channel through which competitive balance would shift.

\subsection{Aggregate Assessment}

The welfare ledger, taken as a whole, is modestly positive. Consumers
received a more dramatic product---more late goals, longer periods of
competitive play---at no apparent cost in total match quality. Players
worked slightly more per match but did not experience higher injury
rates, and the per-minute intensity of play declined. Competitive balance
was unaffected. Broadcasters gained additional minutes of live content.
No constituency for which we have data was made measurably worse off, and
at least one---fans with a taste for late drama---appears to have
benefited. The directive thus represents a rare case of a regulatory
intervention that achieved its stated objective (more effective playing
time) without generating the adverse unintended consequences that so often
accompany well-intentioned rule changes.


%% ---------------------------------------------------------

\section{Conclusion}\label{sec:conclusion}

This paper has studied a natural experiment in time regulation: FIFA's
2022 directive instructing referees to add substantially more stoppage
time to professional football matches. The directive expanded effective
playing time by roughly seven minutes per match, a large and sudden
change to the temporal structure of a globally significant economic
activity. Our central finding is that this expansion of the competitive
window did not increase total goal production. Instead, it redistributed
goals from regulation time into added time while reducing the per-minute
scoring rate, yielding an almost exact offset between the quantity and
intensity margins.

The quantity-intensity tradeoff constitutes the paper's core contribution.
The directive created more competitive time---minutes in which play was
active, contested, and consequential---but the marginal productivity of
each additional minute was lower than the inframarginally displaced
minutes it replaced. Per-minute scoring rates fell by 31 percent in added
time relative to the pre-directive baseline. The result is a
diminishing-returns relationship between allocated time and competitive
output that echoes, in a strikingly different setting, the labor-economics
findings of \citet{Pencavel2015} on the concavity of the hours-output
relationship in manufacturing. When the clock is extended, participants do
not simply produce proportionally more; they adjust effort, strategy, and
intensity in ways that partially---in this case, almost
completely---offset the additional time. The parallel extends to the
mechanism: just as factory workers exhibit declining marginal productivity
in long shifts due to fatigue and pacing, football teams appear to
modulate their competitive intensity when the expected duration of play
increases, conserving effort across a longer horizon.

The finding also connects to the literature on market design and
end-game effects. \citet{Roth2002} demonstrate that hard deadlines in
auction-like settings generate strategic sniping behavior concentrated in
the final moments, whereas soft-close mechanisms distribute activity more
evenly. FIFA's directive functions as a partial soft-close: by
credibly extending the final phase of play, it reduced the scarcity
of late-game minutes and thereby attenuated the strategic premium on
end-of-match actions. The result was a decompression of competitive
activity across a wider temporal window---precisely the outcome that
soft-close mechanisms are designed to produce in auction settings. That
this mechanism operates in a live competitive-sports environment,
without any explicit market-design intent, suggests that the
soft-close principle has broader applicability than its original auction
context.

Several important limitations circumscribe our causal claims. Our
identification relies on the institutional break created by the directive
and the differential exposure of leagues that adopted it at different
times. While the DiD estimates with all controls yield $\beta = 0.744$
(wild-cluster-bootstrap $p = 0.036$) and the permutation test yields $p =
0.045$, we cannot fully eliminate concerns about pre-trends (the joint
$F$-test rejects the null of flat pre-trends at $F = 9.0$, $p < 0.001$).
We interpret our estimates as suggestive of a causal relationship,
supported by the institutional logic of the directive and the consistency
of results across specifications, but we do not claim the identification
is as clean as a randomized experiment. The betting-market evidence
provides complementary support: markets showed no deterioration in
forecast accuracy and no measurable adjustment period, consistent with a
regime change that altered \emph{when} goals occurred without altering
\emph{how many}.

The policy implications extend beyond football. For FIFA, the evidence
suggests that the directive achieved its stated goal of increasing
effective playing time without the adverse consequences that critics
feared: no increase in total goals that might have disrupted the sport's
low-scoring character, no deterioration in competitive balance, and no
measurable increase in player injury rates. For other sports
contemplating similar interventions---extensions of playing time, changes
to clock-management rules, or modifications to end-game
procedures---our results imply that competitive participants will
partially self-regulate in response to expanded time horizons through
intensity adjustment. This self-regulation is not guaranteed to produce a
perfect offset in every setting, but its presence here suggests that
regulators may have more latitude to modify temporal rules than naive
extrapolation would imply. The broader lesson for regulatory economics is
that interventions in competitive environments may be absorbed more
smoothly than anticipated, precisely because strategic agents adjust their
behavior along margins that the regulator did not target.

Finally, our analysis points to a natural extension. Beginning in the
2024--25 season, the English Premier League partially reversed the
directive's emphasis on extended stoppage time, instructing referees to
curtail excessive additions. This reversal provides a mirror-image
natural experiment---a contraction of the competitive window following an
expansion---that will allow future researchers to test whether the
quantity-intensity tradeoff operates symmetrically. If it does, we should
observe a corresponding increase in per-minute intensity as the available
time shrinks, with total goals again remaining stable. The football pitch,
it seems, will continue to serve as an unusually well-instrumented
laboratory for studying how economic agents respond to changes in the
rules of the game.
